% ---------- Множества ----------
\newcommand{\R}{\mathbb{R}}
\newcommand{\N}{\mathbb{N}}
\newcommand{\Z}{\mathbb{Z}}

% ---------- Вероятность и моменты ----------
\renewcommand{\P}{\mathbb{P}}                       % вероятность
\newcommand{\E}{\mathbb{E}}                         % матожидание
\DeclareMathOperator{\D}{\mathbb{D}}                % дисперсия
\DeclareMathOperator{\Var}{Var}
\DeclareMathOperator{\cov}{cov}
\DeclareMathOperator{\corr}{corr}
\newcommand{\Ind}[1]{\mathbb{I}\{#1\}}              % индикатор события

% ---------- Параметрическая модель ----------
\newcommand{\Pt}{\P_\theta}
\newcommand{\Et}{\E_\theta}
\newcommand{\Dt}{\D_\theta}
\newcommand{\Th}{\Theta}                            % параметрическое множество
\newcommand{\sample}{X_1, \ldots, X_n}              % выборка
\newcommand{\Xbar}{\overline{X}}                    % выборочное среднее
\newcommand{\est}[1]{\widehat{#1}}                  % оценка
\newcommand{\Pemp}{\P_n^{*}}                        % эмпирическое распределение
\newcommand{\iid}{\overset{\textup{i.i.d.}}{\sim}}
\DeclareMathOperator*{\argmax}{arg\,max}
\DeclareMathOperator*{\argmin}{arg\,min}

% ---------- Сходимости ----------
\newcommand{\toas}{\xrightarrow{\textup{п.н.}}}
\newcommand{\toP}{\xrightarrow{\P}}
\newcommand{\tod}{\xrightarrow{d}}
\newcommand{\toLp}[1]{\xrightarrow{L_{#1}}}

% ---------- Распределения ----------
\newcommand{\Norm}{\mathcal{N}}
\newcommand{\Bern}{\mathrm{Bern}}
\newcommand{\Bin}{\mathrm{Bin}}
\newcommand{\Pois}{\mathrm{Pois}}
\newcommand{\Exp}{\mathrm{Exp}}
\newcommand{\Unif}{\mathrm{U}}

% ---------- Прочее ----------
% Наклонные знаки ≤ и ≥ во всём конспекте (как принято в русских учебниках)
\renewcommand{\le}{\leqslant}
\renewcommand{\leq}{\leqslant}
\renewcommand{\ge}{\geqslant}
\renewcommand{\geq}{\geqslant}
\newcommand{\eqdef}{\coloneqq}
\newcommand{\abs}[1]{\lvert #1 \rvert}
\renewcommand{\epsilon}{\varepsilon}
\renewcommand{\phi}{\varphi}
